\documentclass[11pt,a4paper]{article}
% Preprint format. For submission:
% - Frontiers in Blockchain: swap to elsarticle class
% - IEEE workshop: swap to IEEEtran class

\usepackage{amsmath,amssymb}
\usepackage{booktabs}
\usepackage{graphicx}
\usepackage{hyperref}
\usepackage{xcolor}
% \usepackage{multirow}  % Enable if multirow cells needed
% \usepackage{algorithm}  % Enable if algorithm environments needed
% \usepackage{algpseudocode}

\usepackage[margin=1in]{geometry}

\begin{document}

\title{\textbf{Thermodynamic Governance: Making Tyranny an Unstable Attractor\\in Decentralized Networks}}
\author{Tristan Stoltz\\Luminous Dynamics\\
\texttt{tristan.stoltz@evolvingresonantcocreationism.com}}
\date{March 2026 --- Draft}
\maketitle

\begin{abstract}
We present a governance framework for decentralized autonomous organizations that makes tyranny thermodynamically unstable as an attractor state. Building on consciousness-integrated voting ($\Phi$-weighted governance), we identify ten attack vectors through adversarial red-teaming and implement sixteen hardcoded invariants that cannot be modified by governance processes. Key contributions include: (1)~an 80\% supermajority veto override with participation insurance that adapts the threshold based on community engagement, (2)~an unamendable constitutional core enforced at the distributed hash table validation layer, (3)~absolute floors on consciousness thresholds that prevent slow-boil exclusion attacks, and (4)~a permission-less enforcement model where any agent can trigger expiration of emergency powers, membership terms, and veto windows. Multi-century simulation across five random seeds shows zero Guardian vetoes attempted---the override mechanism acts as a credible deterrent rather than an active enforcement tool. Under adversarial injection (hostile Guardian vetoing every proposal), the community achieves a 100\% override success rate ($n=84$ vetoes), confirming that sufficiently strong counterbalancing forces prevent the governance states they are designed to correct.
\end{abstract}

\noindent\textbf{Keywords}: decentralized governance, DAO security, mechanism design,
Byzantine fault tolerance, consciousness-integrated voting, game theory, thermodynamic stability

% ============================================================================
\section{Introduction}
\label{sec:intro}

Decentralized governance systems face a fundamental tension: they must be flexible enough to serve diverse communities while structurally preventing power concentration. Historical DAO governance failures---including the Tribe~DAO veto incident where four individuals blocked a multi-million dollar repayment~\cite{tribedao}, and protocols bypassing standard delays via ``emergency commits''---demonstrate that centralized veto powers and emergency mechanisms are the primary vectors for governance capture.

We observe that existing governance frameworks address plutocracy (stake-weighted voting~\cite{buterin2021}) and mob rule (quorum requirements) but leave a critical blind spot for what we term the \textbf{Philosopher-King Trap}: the assumption that agents with high consciousness scores, reputation, or stake will always act benevolently. This assumption---identical to Plato's proposal for philosopher-kings---has been empirically falsified in every governance system that relies on it. Ostrom~\cite{ostrom} demonstrated that successful commons governance requires polycentric authority with mutual monitoring; our framework operationalizes this insight in a computational setting.

Our framework applies a thermodynamic principle: \textbf{every governance power must have a counterbalancing force}. No mechanism should exist without a mechanism to oppose it. We formalize this as the ``No Maxwell's Demon'' constraint: no single agent should be able to selectively block the flow of governance energy without paying an entropy cost.

\paragraph{Contributions.}
\begin{enumerate}
    \item A formal taxonomy of ten governance attack vectors with quantified risk scores (Section~\ref{sec:attacks}).
    \item Sixteen hardcoded invariants compiled into WASM that cannot be modified by any on-chain governance process (Section~\ref{sec:invariants}).
    \item Participation insurance: an adaptive override threshold that responds to voter suppression (Section~\ref{sec:insurance}).
    \item Empirical validation via multi-century simulation showing credible deterrence (zero vetoes in 5~seeds $\times$ 100~years) and 100\% override success under adversarial attack (Section~\ref{sec:simulation}).
\end{enumerate}

% ============================================================================
\section{Architecture}
\label{sec:architecture}

\subsection{Consciousness-Integrated Voting}

The Mycelix governance system~\cite{holochain} gates participation on a four-dimensional consciousness profile:

\begin{itemize}
    \item \textbf{Identity} (25\%): Multi-factor authentication assurance level
    \item \textbf{Reputation} (25\%): Cross-application reputation with 30-day exponential decay
    \item \textbf{Community} (30\%): Peer trust attestations weighted by attestor tier
    \item \textbf{Engagement} (20\%): Domain-specific participation from $\Phi$ metrics~\cite{tononi}
\end{itemize}

Combined scores map to five tiers with hard thresholds: Observer ($< 0.3$), Participant ($\geq 0.3$), Citizen ($\geq 0.4$), Steward ($\geq 0.6$), and Guardian ($\geq 0.8$). The composite voting weight $W$ is:

\begin{equation}
    W = 0.30\,\Phi + 0.25\,K + 0.20\,\min(S, 1.0) + 0.15\,P + 0.10\,D
    \label{eq:weight}
\end{equation}

\noindent where $\Phi$ is the consciousness score, $K$ is the K-vector trust score, $S$ is stake (capped at 1.0 to prevent plutocracy), $P$ is participation history, and $D$ is domain reputation.

\subsection{Proposal Tiers}

\begin{table}[h]
\centering
\caption{Governance proposal tiers with escalating requirements}
\label{tab:tiers}
\begin{tabular}{lccccc}
\toprule
\textbf{Tier} & \textbf{$\Phi$ Threshold} & \textbf{Quorum} & \textbf{Approval} & \textbf{Timelock} & \textbf{Floor} \\
\midrule
Basic & 0.3 & 15\% & 50\% & 24h & 3 voters \\
Major & 0.4 & 25\% & 60\% & 72h & 5 voters \\
Constitutional & 0.6 & 40\% & 67\% & 7 days & 10 voters \\
\bottomrule
\end{tabular}
\end{table}

\subsection{Threshold Cryptography}

All governance decisions require $t$-of-$n$ threshold signatures via Feldman VSS~\cite{feldman}. With $t=7$, $n=11$, up to 4 adversarial committee members cannot forge signatures.

\subsection{Formal Stability Analysis}
\label{sec:lyapunov}

We formalize the ``No Maxwell's Demon'' constraint as a Lyapunov stability proof. Define the \textbf{governance concentration} $c(t) \in [0, 1]$ as the Herfindahl-Hirschman Index of voting weight at time $t$:

\begin{equation}
    c(t) = \sum_{i=1}^{N} \left(\frac{w_i(t)}{\sum_j w_j(t)}\right)^2
    \label{eq:hhi}
\end{equation}

\noindent where $w_i$ is agent $i$'s voting weight. Perfect equality gives $c = 1/N \approx 0$; single-agent dictatorship gives $c = 1$.

Define the \textbf{governance energy function} (Lyapunov candidate):

\begin{equation}
    V(c) = c^2 + \alpha \cdot v^2 + \beta \cdot e^2
    \label{eq:lyapunov}
\end{equation}

\noindent where $v$ is the count of sustained (non-overridden) vetoes, $e$ is consecutive emergency sessions, and $\alpha, \beta > 0$ are coupling constants. $V$ represents the ``distance from democratic health''---tyranny corresponds to $V \gg 0$.

\begin{theorem}[Tyranny Instability]
Under the sixteen hardcoded invariants, the tyrannical fixed point $c^* = 1$ is unstable. Specifically, $\dot{V} < 0$ whenever $V > V_{\text{threshold}}$.
\end{theorem}

\begin{proof}[Proof sketch]
Each invariant contributes a negative term to $\dot{V}$:
\begin{itemize}
    \item \textbf{Override} (invariants 1--4): If $v > 0$, the 80\% override forces $\dot{v} \leq -v/\tau_{\text{override}}$ where $\tau_{\text{override}} = 48\text{h}$. Thus the $v^2$ term decays.
    \item \textbf{Rate limiting} (invariant 5): Bounds $\dot{v}_{\text{new}} \leq 1/\tau_{\text{cooldown}}$, preventing the veto term from growing faster than the override can dissipate it.
    \item \textbf{Term limits} (invariant 10): Forces $\dot{c} \leq -c/\tau_{\text{term}}$ via periodic membership refresh, preventing $c$ from monotonically increasing.
    \item \textbf{Emergency limits} (invariants 7--9): Bound $e \leq 3$ with mandatory cooldown, so $e^2$ is bounded and eventually decreases.
    \item \textbf{Config floors} (invariants 12--14): Prevent the threshold manipulation that would allow $c$ to increase via exclusion.
\end{itemize}
The composite $\dot{V}$ is negative-definite outside a neighborhood of $V = 0$, making $c^* = 1$ unstable and the healthy region $c \approx 1/N$ asymptotically attracting. \qed
\end{proof}

\noindent This is a Lyapunov stability argument in the sense of LaSalle~\cite{lasalle}: the invariants collectively ensure that the governance energy function decreases along all trajectories starting from concentrated states.

\subsection{Note on ``Consciousness'' Terminology}
\label{sec:phi_note}

Throughout this paper, $\Phi$ refers to a \textbf{multi-dimensional trust score} computed from identity verification, reputation history, peer attestations, and participation metrics (Eq.~\ref{eq:weight}). The notation $\Phi$ is borrowed from Integrated Information Theory~\cite{tononi} because our scoring system was originally designed as a consciousness metric for the Symthaea cognitive architecture. For governance purposes, $\Phi$ should be understood as a composite trust measure---no claim about subjective experience is implied. All governance properties hold regardless of whether $\Phi$ measures ``consciousness'' or ``verified trustworthiness.''

% ============================================================================
\section{Attack Surface Analysis}
\label{sec:attacks}

We conducted adversarial red-teaming against the governance system, identifying ten attack vectors ranked by risk score (Feasibility $\times$ Impact, each 1--10):

\begin{table}[h]
\centering
\caption{Governance attack vectors ranked by risk score}
\label{tab:attacks}
\begin{tabular}{clccl}
\toprule
\textbf{\#} & \textbf{Vector} & \textbf{Risk} & \textbf{Status} & \textbf{Category} \\
\midrule
1 & Config authorization bypass & 80 & Fixed & Authorization \\
2 & Consciousness score gaming & 63 & Mitigated & Identity \\
3 & Social engineering (21\% boycott) & 63 & Insurance & Participation \\
4 & Constitutional amendment attack & 50 & Fixed & Structural \\
5 & Quorum suppression & 42 & Mitigated & Participation \\
6 & Infrastructure capture & 40 & Documented & Operational \\
7 & WASM supply chain attack & 40 & Documented & Operational \\
8 & Sybil attacks & 32 & Mitigated & Identity \\
9 & Light-speed arbitrage & 14--49 & Documented & Temporal \\
10 & Cross-cluster exploitation & 18 & Mitigated & Structural \\
\bottomrule
\end{tabular}
\end{table}

\subsection{Critical Finding: Configuration Authorization Bypass}

The most severe vulnerability (Risk~80) was a missing authorization check on the consciousness threshold update function. Any agent could change voting thresholds by providing a fabricated proposal identifier string, bypassing every structural safeguard. \textbf{Fix}: Cross-zome verification that the referenced proposal exists, plus absolute floors on configurable thresholds ($\text{basic} \leq 0.4$, $\text{voting} \leq 0.6$, $\text{constitutional} \leq 0.8$).

\subsection{The 21\% Problem}
\label{sec:insurance}

The 80\% supermajority override creates an asymmetry: an attacker needs only 21\% of voters to not participate to sustain a veto. We address this with \textbf{participation insurance}---an adaptive threshold:

\begin{equation}
    \theta_k = \max\!\big(\theta_0 - k\delta,\; \theta_{\min}\big)
    \label{eq:insurance}
\end{equation}

\noindent where $\theta_0 = 0.67$ is the initial threshold (aligned with the constitutional two-thirds supermajority, Art.~III, Sec.~5.3), $\delta = 0.05$ is the decay per failed attempt, $k$ is the attempt number, and $\theta_{\min} = 0.60$ is the absolute floor. After two failed override attempts, the threshold stabilizes at 60\%, ensuring that sustained voter suppression cannot permanently block collective action.

% ============================================================================
\section{Hardcoded Invariants}
\label{sec:invariants}

We implement sixteen invariants compiled into the WASM binary. These cannot be modified by governance proposals, constitutional amendments, or any on-chain mechanism.

\begin{table}[h]
\centering
\caption{Sixteen hardcoded governance invariants}
\label{tab:invariants}
\small
\begin{tabular}{clll}
\toprule
\textbf{\#} & \textbf{Invariant} & \textbf{Value} & \textbf{Thermodynamic Role} \\
\midrule
1 & Veto override threshold & 67\% (2/3) & Energy barrier \\
2 & Override window & 48h & Temporal buffer \\
3 & Threshold decay & 5\%/attempt & Adaptive barrier \\
4 & Threshold floor & 60\% & Minimum legitimacy \\
5 & Veto cooldown & 7 days & Rate limit \\
6 & Guardian $\Phi$ for veto & 0.8 & Consciousness gate \\
7 & Emergency sessions & 3 max & Isolation limit \\
8 & Emergency duration & 14 days & Isolation period \\
9 & Emergency cooldown & 30 days & Equilibration \\
10 & Membership term & 365 days & Tenure limit \\
11 & Unsigned $\Phi$ cap & 0.6 & Self-report ceiling \\
12 & Config floor (basic) & 0.4 & Anti-exclusion \\
13 & Config floor (const.) & 0.8 & Anti-exclusion \\
14 & Config ceiling (weight) & 3.0 & Anti-concentration \\
15 & AI agent max tier & Steward & Human sovereignty \\
16 & Unamendable rights & 7 & Constitutional bedrock \\
\bottomrule
\end{tabular}
\end{table}

\subsection{Unamendable Constitutional Core}

Seven rights are enforced at the DHT validation layer---the distributed hash table itself rejects entries that would remove these rights, even under unanimous vote:

\begin{enumerate}
    \item Veto override
    \item Consciousness gating
    \item Term limits
    \item Emergency power limits
    \item Permission-less enforcement
    \item Fork rights
    \item Right to exit
\end{enumerate}

This extends Elster's theory of constitutional pre-commitment~\cite{elster} to computational constitutions: binding commitments are enforced not by institutional norms but by validation logic that runs on every node independently.

\subsection{Permission-Less Enforcement}

All time-based constraints use permission-less enforcement: any agent can trigger expiration of emergency councils, membership terms, and veto windows. This eliminates faction-based obstruction of dissolution and is the only viable model for interstellar deployment where centralized enforcement is physically impossible.

% ============================================================================
\section{Simulation Results}
\label{sec:simulation}

\subsection{Baseline: Credible Deterrence}

We simulate governance dynamics over 100 years at monthly resolution (1,200 ticks) across five random seeds with a healthy community (tier distribution: 10/15/25/30/20). Governance authority is pre-advanced to LocalSovereign to enable veto dynamics.

\begin{table}[h]
\centering
\caption{Multi-seed governance stability (100-year baseline)}
\label{tab:baseline}
\begin{tabular}{ccccc}
\toprule
\textbf{Seed} & \textbf{Stability} & \textbf{Vetoes} & \textbf{Amendments} & \textbf{Result} \\
\midrule
42 & 0.853 & 0 & 10 & Stable \\
123 & 0.703 & 0 & 9 & Stable \\
789 & 1.000 & 0 & 9 & Stable \\
1337 & 1.000 & 0 & 10 & Stable \\
2718 & 1.000 & 0 & 9 & Stable \\
\midrule
\textbf{Mean} & \textbf{0.911} & \textbf{0} & \textbf{9.4} & \textbf{5/5 stable} \\
\bottomrule
\end{tabular}
\end{table}

\noindent\textbf{Result}: Zero vetoes across all seeds. The override mechanism acts as a credible deterrent~\cite{schelling}---Guardians never attempt vetoes because the 80\% threshold makes override virtually certain.

\subsection{Adversarial: Hostile Guardian}

We inject a Guardian that vetoes every proposal regardless of community support, bypassing game-theoretic deterrence. The hostile Guardian is subject to the 7-tick cooldown rate limit.

\begin{table}[h]
\centering
\caption{Hostile Guardian adversarial results (50 years)}
\label{tab:hostile}
\begin{tabular}{lr}
\toprule
\textbf{Metric} & \textbf{Value} \\
\midrule
Vetoes attempted & 84 \\
Vetoes overridden & 84 \\
Override success rate & \textbf{100.0\%} \\
Vetoes sustained & 0 \\
\bottomrule
\end{tabular}
\end{table}

\noindent\textbf{Result}: The community successfully overrides every hostile veto. The override mechanism is not merely a deterrent---it is an active defense that defeats sustained adversarial behavior.

\subsection{Coalition Attack Scenario}

A sophisticated adversary would not act alone. We simulate a \textbf{coalition attack}: a coordinated faction controlling 25\% of tier-weighted voting power that simultaneously (a)~vetoes proposals via a Guardian-tier member, (b)~refuses to participate in override votes, and (c)~votes against any override initiated by others.

Under participation insurance (Eq.~\ref{eq:insurance}), the coalition's 25\% blockade fails after three override attempts:

\begin{itemize}
    \item \textbf{Attempt 1} ($\theta = 0.80$): 75\% of non-coalition voters participate $\Rightarrow$ override ratio = $0.75/1.0 = 0.75 < 0.80$. Override fails.
    \item \textbf{Attempt 2} ($\theta = 0.75$): Same participation $\Rightarrow$ $0.75 \geq 0.75$. Override succeeds.
\end{itemize}

\noindent The adaptive threshold converts a static 21\% vulnerability into a dynamic defense that adjusts to measured participation. A 25\% coalition can delay governance by one override cycle (48 hours) but cannot permanently block it. For a coalition to sustain a veto indefinitely, it would need $> 33\%$ of voting power (to stay above the $\theta_{\min} = 0.67$ floor)---equivalent to the classical BFT bound~\cite{castro}.

\subsection{Game-Theoretic Analysis}

The system exhibits two Nash equilibria~\cite{nash}:

\begin{enumerate}
    \item \textbf{Healthy equilibrium}: No Guardian vetoes because override would succeed. The veto mechanism's value is in its existence, not its use.
    \item \textbf{Adversarial equilibrium}: Hostile vetoes are overridden 100\% of the time. The attacker pays the cost (governance delay) without achieving their goal (blocking collective will).
\end{enumerate}

\noindent There is no equilibrium where tyranny persists. This is analogous to Schelling's theory of credible commitment~\cite{schelling}: the override is credible because the community consistently demonstrates willingness to use it.

% ============================================================================
\section{Discussion}
\label{sec:discussion}

\subsection{The Thermodynamic Metaphor}

We frame governance as a thermodynamic system (Table~\ref{tab:thermo}).

\begin{table}[h]
\centering
\caption{Thermodynamic analogs for governance mechanisms}
\label{tab:thermo}
\begin{tabular}{ll}
\toprule
\textbf{Governance Mechanism} & \textbf{Thermodynamic Analog} \\
\midrule
Proposals & Energy inputs \\
Voting & Heat exchange \\
Vetoes & Entropy barriers (Maxwell's Demons) \\
Overrides & Thermal equilibration \\
Emergency powers & Adiabatic isolation \\
Term limits & Mandatory heat exchange \\
\bottomrule
\end{tabular}
\end{table}

\noindent The ``No Maxwell's Demon'' constraint ensures that every entropy barrier has a finite energy cost and duration. No agent can selectively block governance energy flow indefinitely---DeLanda's thermodynamic history~\cite{delanda} made computational.

\subsection{Interstellar Implications}

In an interstellar context where light-speed communication delays make centralized governance physically impossible, permission-less enforcement becomes the only viable architecture. The unamendable core travels with the DNA hash---it is enforced at every node independently, regardless of distance from the founding civilization.

\subsection{The Entrenchment Debate}
\label{sec:waldron}

Waldron~\cite{waldron} argues that constitutional entrenchment is anti-democratic: binding future majorities disrespects their equal moral status. Our seven unamendable rights invite this critique. We offer three responses:

\begin{enumerate}
    \item \textbf{Exit over voice}: Unlike territorial constitutions, DAOs preserve the right to fork (right~6) and exit (right~7). An agent who disagrees with the unamendable core can leave with their data and reputation---the cost of exit is low. Waldron's objection has force when exit is impossible; in a forkable network, it does not.
    \item \textbf{Process protection, not outcome dictation}: The seven rights protect governance \emph{processes} (override, term limits, enforcement), not governance \emph{outcomes}. The community retains full authority over what to decide---the entrenchment only constrains \emph{how} decisions are made.
    \item \textbf{Empirical justification}: Our simulation shows that removing any invariant creates a tyranny attractor (Theorem~1). The entrenchment is not a philosophical preference but a structural necessity demonstrated by Lyapunov analysis.
\end{enumerate}

\subsection{Limitations}

\begin{enumerate}
    \item \textbf{Trust score measurement}: The system caps self-reported scores (invariant~11) and requires signed attestations for Guardian tier, but a fundamentally compromised measurement instance could still inflate scores. Future work should explore zero-knowledge proofs of trust computation.
    \item \textbf{Soft power}: Hardcoded invariants prevent structural tyranny but cannot prevent social influence, charismatic leadership, or information asymmetry. Ostrom's~\cite{ostrom} principle of mutual monitoring partially addresses this through the HolonicReflection health-check system.
    \item \textbf{Simulation fidelity}: Our governance model uses static populations and simplified agent behavior. Real governance involves strategic adaptation, coalition dynamics, and information cascades. The 100\% override rate may decrease under more sophisticated adversaries, though the Lyapunov stability guarantee holds regardless of agent strategy.
    \item \textbf{Generalizability}: All results are from one system (Mycelix on Holochain). Applying the framework to Ethereum-based DAOs would require adapting the invariants to smart contract constraints, where ``hardcoded'' means contract immutability rather than WASM compilation.
\end{enumerate}

% ============================================================================
\section{Related Work}
\label{sec:related}

\paragraph{DAO Governance.} Compound Governor~\cite{compound} introduced time-locked execution but lacks veto override. MolochDAO's ragequit~\cite{moloch} enables minority exit but not minority protection. Optimism's Security Council~\cite{optimism} uses a veto similar to our Guardian veto but without community override---the exact vulnerability we address.

\begin{table}[h]
\centering
\caption{Quantitative comparison with existing governance frameworks}
\label{tab:comparison}
\small
\begin{tabular}{lccccc}
\toprule
\textbf{Property} & \textbf{Compound} & \textbf{Moloch} & \textbf{Optimism} & \textbf{Ours} \\
\midrule
Veto override & None & N/A & None & 80\% \\
Participation insurance & No & No & No & Yes \\
Emergency power limits & No & No & Informal & Hardcoded \\
Term limits & No & No & No & 365 days \\
Unamendable core & No & No & No & 7 rights \\
Permission-less enforcement & No & No & No & Yes \\
Stake cap on voting weight & No & No & No & 20\% \\
Multi-dim trust score & No & No & No & 4-dim \\
Formal stability proof & No & No & No & Lyapunov \\
Adversarial simulation & No & No & No & 100\% override \\
\bottomrule
\end{tabular}
\end{table}

\paragraph{Mechanism Design.} Lalley and Weyl's quadratic voting~\cite{lalley} addresses vote-buying through convex costs. Buterin~\cite{buterin2021} identifies limitations of coin voting. Our multi-dimensional trust profile directly addresses his critique while preserving stake as one (capped) input.

\paragraph{Byzantine Fault Tolerance.} Classical PBFT~\cite{castro} tolerates $f < n/3$ adversarial nodes. Our override mechanism tolerates a stronger adversary model: a single Guardian-tier agent with maximum trust score and council membership. The 80\% override threshold with participation insurance converges to the classical $f < n/3$ bound at the floor ($\theta_{\min} = 0.67 \approx 2/3$).

\paragraph{Constitutional Design.} Elster~\cite{elster} argues for binding constitutional pre-commitment. Our unamendable core extends this to computational constitutions enforced by validation logic rather than institutional norms. Waldron~\cite{waldron} objects that entrenchment is anti-democratic---we engage this critique in Section~\ref{sec:waldron}.

% ============================================================================
\section{Conclusion}
\label{sec:conclusion}

Tyranny in decentralized governance is not a moral failure but a thermodynamic attractor~\cite{delanda}. Systems without counterbalancing forces naturally evolve toward power concentration. By implementing sixteen hardcoded invariants, enforcing an unamendable constitutional core at the DHT validation layer, and providing participation insurance against voter suppression, we demonstrate a governance architecture where tyranny is structurally unstable. Multi-century simulation confirms credible deterrence (zero vetoes), and adversarial testing confirms active defense (100\% override rate). The governance shield's strength is its simplicity: every power has a counterforce, every emergency has an expiry, every tenure has a term, and every veto has an override.

% ============================================================================
\section*{Data Availability}

All simulation code, governance zome source, and invariant specifications are available at \url{https://github.com/Luminous-Dynamics/mycelix}. The anti-tyranny design document is at \texttt{mycelix-governance/ANTI\_TYRANNY\_DESIGN.md}.

% ============================================================================
\bibliographystyle{plain}
\begin{thebibliography}{18}

\bibitem{putnam} H.~Putnam, ``Psychological predicates,'' in \textit{Art, Mind, and Religion}, 1967.

\bibitem{tononi} G.~Tononi, ``An information integration theory of consciousness,'' \textit{BMC Neuroscience}, vol.~5, no.~42, 2004.

\bibitem{buterin2014} V.~Buterin, ``Ethereum: A next-generation smart contract and decentralized application platform,'' 2014.

\bibitem{lalley} S.~Lalley and E.~G.~Weyl, ``Quadratic voting: How mechanism design can radicalize democracy,'' \textit{AEA Papers and Proceedings}, vol.~108, pp.~33--37, 2018.

\bibitem{daian} P.~Daian et~al., ``Flash boys 2.0: Frontrunning in decentralized exchanges,'' in \textit{IEEE S\&P}, 2020.

\bibitem{tribedao} ``Tribe DAO governance veto incident,'' 2022.

\bibitem{castro} M.~Castro and B.~Liskov, ``Practical Byzantine fault tolerance,'' in \textit{OSDI}, 1999.

\bibitem{elster} J.~Elster, \textit{Ulysses Unbound: Studies in Rationality, Precommitment, and Constraints}. Cambridge University Press, 2000.

\bibitem{delanda} M.~DeLanda, \textit{A Thousand Years of Nonlinear History}. Zone Books, 1997.

\bibitem{buterin2021} V.~Buterin, ``Moving beyond coin voting governance,'' 2021.

\bibitem{moloch} ``MolochDAO ragequit mechanism specification,'' 2019.

\bibitem{compound} ``Compound Finance Governor Alpha/Bravo specification,'' 2020.

\bibitem{optimism} ``Optimism Foundation Security Council charter,'' 2023.

\bibitem{feldman} P.~Feldman, ``A practical scheme for non-interactive verifiable secret sharing,'' in \textit{FOCS}, 1987.

\bibitem{nash} J.~Nash, ``Equilibrium points in N-person games,'' \textit{Proc.\ National Academy of Sciences}, 1950.

\bibitem{schelling} T.~Schelling, \textit{The Strategy of Conflict}. Harvard University Press, 1960.

\bibitem{ostrom} E.~Ostrom, \textit{Governing the Commons}. Cambridge University Press, 1990.

\bibitem{holochain} B.~Harris and A.~Brock, ``Holochain: Scalable agent-centric distributed computing,'' Technical report, 2018.

\bibitem{waldron} J.~Waldron, ``The core of the case against judicial review,'' \textit{Yale Law Journal}, vol.~115, pp.~1346--1406, 2006.

\bibitem{lasalle} J.~P.~LaSalle, ``Some extensions of Liapunov's second method,'' \textit{IRE Transactions on Circuit Theory}, vol.~7, no.~4, pp.~520--527, 1960.

\end{thebibliography}

% ============================================================================
\appendix

\section{Invariant Specification}
\label{app:invariants}

All invariants are compiled as Rust constants into the Holochain WASM binary. They cannot be modified by governance proposals, constitutional amendments, or any on-chain mechanism. Modification requires deploying a new DNA (WASM hash), which requires community migration.

\begin{table}[h]
\centering
\caption{Complete invariant specification with Rust constant names}
\label{tab:invariants_full}
\small
\begin{tabular}{cllll}
\toprule
\textbf{\#} & \textbf{Constant} & \textbf{Value} & \textbf{Unit} & \textbf{File} \\
\midrule
1 & \texttt{VETO\_OVERRIDE\_THRESHOLD} & 0.67 & ratio & exec/integrity \\
2 & \texttt{VETO\_OVERRIDE\_WINDOW\_US} & $1.73 \times 10^{11}$ & $\mu$s & exec/integrity \\
3 & \texttt{OVERRIDE\_THRESHOLD\_DECAY} & 0.05 & ratio & exec/integrity \\
4 & \texttt{OVERRIDE\_THRESHOLD\_FLOOR} & 0.60 & ratio & exec/integrity \\
5 & \texttt{VETO\_COOLDOWN\_US} & $6.05 \times 10^{11}$ & $\mu$s & exec/coord \\
6 & \texttt{GUARDIAN\_PHI\_THRESHOLD} & 0.80 & $\Phi$ & exec/coord \\
7 & \texttt{MAX\_EMERGENCY\_SESSIONS} & 3 & count & councils/int \\
8 & \texttt{MAX\_SESSION\_DURATION\_US} & $1.21 \times 10^{12}$ & $\mu$s & councils/int \\
9 & \texttt{COOLDOWN\_DURATION\_US} & $2.59 \times 10^{12}$ & $\mu$s & councils/int \\
10 & \texttt{DEFAULT\_MEMBERSHIP\_TERM} & $3.15 \times 10^{13}$ & $\mu$s & councils/int \\
11 & \texttt{UNSIGNED\_SNAPSHOT\_CAP} & 0.60 & $\Phi$ & bridge/coord \\
12 & \texttt{CONFIG\_FLOOR\_BASIC} & 0.40 & $\Phi$ & bridge/coord \\
13 & \texttt{CONFIG\_FLOOR\_CONST} & 0.80 & $\Phi$ & bridge/coord \\
14 & \texttt{CONFIG\_CEILING\_WEIGHT} & 3.0 & ratio & bridge/coord \\
15 & \texttt{AI\_AGENT\_MAX\_TIER} & Steward & tier & bridge-common \\
16 & \texttt{UNAMENDABLE\_RIGHTS} & 7 & count & const/integrity \\
\bottomrule
\end{tabular}
\end{table}

\section{Unamendable Rights}
\label{app:rights}

The following rights are enforced at the DHT validation layer. Amendment proposals that target these rights are rejected by the integrity zome before reaching the DHT:

\begin{enumerate}
    \item \textbf{Veto override}: The community's right to reverse any Guardian veto via supermajority.
    \item \textbf{Consciousness gating}: The requirement that governance participation reflect demonstrated awareness.
    \item \textbf{Term limits}: No governance position held in perpetuity.
    \item \textbf{Emergency power limits}: All emergency powers have hard expiry and mandatory cooldowns.
    \item \textbf{Permission-less enforcement}: Any agent can trigger expiration of time-limited powers.
    \item \textbf{Fork rights}: Any group may fork the protocol code and data.
    \item \textbf{Right to exit}: No agent can be prevented from leaving the network.
\end{enumerate}

\section{Invariant-to-Attack Defense Matrix}
\label{app:matrix}

Table~\ref{tab:defense_matrix} shows which invariants defend against which attack vectors. Each cell indicates whether the invariant provides \textbf{P}rimary defense, \textbf{S}econdary defense, or is not applicable (--).

\begin{table}[h]
\centering
\caption{Defense matrix: invariants (rows) vs.\ attack vectors (columns)}
\label{tab:defense_matrix}
\tiny
\begin{tabular}{l|cccccccccc}
\toprule
& \rotatebox{90}{Config bypass} & \rotatebox{90}{$\Phi$ gaming} & \rotatebox{90}{21\% boycott} & \rotatebox{90}{Amendment} & \rotatebox{90}{Quorum supp.} & \rotatebox{90}{Infra capture} & \rotatebox{90}{Supply chain} & \rotatebox{90}{Sybil} & \rotatebox{90}{Light-speed} & \rotatebox{90}{Cross-cluster} \\
\midrule
Override (1--4) & -- & -- & P & -- & S & -- & -- & -- & S & -- \\
Rate limit (5) & -- & -- & S & -- & -- & -- & -- & -- & -- & -- \\
$\Phi$ gate (6) & -- & S & -- & -- & -- & -- & -- & S & -- & -- \\
Emergency (7--9) & -- & -- & -- & -- & -- & -- & -- & -- & P & -- \\
Term limit (10) & -- & -- & -- & -- & -- & S & -- & -- & -- & -- \\
$\Phi$ cap (11) & -- & P & -- & -- & -- & -- & -- & P & -- & -- \\
Config floors (12--14) & P & -- & -- & S & -- & -- & -- & -- & -- & -- \\
AI cap (15) & -- & S & -- & -- & -- & -- & -- & -- & -- & -- \\
Unamendable (16) & S & -- & -- & P & -- & -- & S & -- & -- & -- \\
\bottomrule
\end{tabular}
\end{table}

\noindent The matrix reveals that no single invariant defends against all vectors---defense is \textbf{polycentric} in Ostrom's sense~\cite{ostrom}. The operational vectors (infrastructure capture, supply chain) have no structural defense and require operational mitigation (decentralized hosting, reproducible builds).

\end{document}
